% BHU PYQ -- Transcribed & Corrected
% Paper: GGRSE21: Basics of Remote Sensing (Mid-Term)
% Exam: B.A. / B.Sc. Semester III Mid-Term Examination 2024-25 (NEP SEC)
% Department: Geography

\documentclass[11pt,a4paper]{article}
\usepackage[a4paper,top=2.5cm,bottom=2.5cm,left=2.8cm,right=2.8cm,headheight=14pt]{geometry}
\usepackage{amsmath,amssymb}
\usepackage[T1]{fontenc}
\usepackage[utf8]{inputenc}
\usepackage{lmodern,microtype,fancyhdr,enumitem,hyperref,lastpage,parskip}

\hypersetup{
    colorlinks=true,
    urlcolor=black,
    linkcolor=black,
    pdftitle={GGRSE21: Basics of Remote Sensing -- Mid-Term 2024-25}
}

\pagestyle{fancy}
\fancyhf{}
\renewcommand{\headrulewidth}{0.4pt}
\renewcommand{\footrulewidth}{0.4pt}
\fancyhead[L]{\small\textit{Banaras Hindu University}}
\fancyhead[R]{\small\textit{GGRSE21: Basics of Remote Sensing (Mid-Term)}}
\fancyfoot[C]{\small Page \thepage\ of \pageref{LastPage}}

\newcommand{\pts}[1]{\hfill\textit{[#1]}}

\begin{document}

\begin{center}
  {\large\bfseries BANARAS HINDU UNIVERSITY}\\[2pt]
  {\normalsize B.A. / B.Sc. Semester III Mid-Term Examination 2024-25 (NEP SEC)}\\[6pt]
  \rule{0.8\linewidth}{0.5pt}\\[6pt]
  {\large\bfseries DEPARTMENT OF GEOGRAPHY}\\[4pt]
  {\large\bfseries Paper Code: GGRSE21 : Basics of Remote Sensing}\\[4pt]
  {\normalsize Time: 1 Hour\hfill Maximum Marks: 20}
\end{center}

\rule{\linewidth}{0.4pt}

\smallskip
\noindent\textit{Note: Answer any two questions. Each question carries 10 marks.\hfill [$2 \times 10 = 20$ Marks]}

\bigskip

\begin{enumerate}[label=\textbf{Q\arabic*.}]
    \item Explain the different types of aerial photographs and their characteristics in detail.\pts{10}
    \item Briefly describe the principles and scope of remote sensing, citing suitable examples.\pts{10}
    \item Describe the difference between airborne and spaceborne remote sensing, citing suitable examples.\pts{10}
\end{enumerate}

\end{document}
